% ID: SEL-D0030-OB-E02
% ESTADO: validada. S031, cotejada con examen y criterios.
\textbf{a)} En $1$ la continuidad exige $1+a=b+2$.
Las derivadas de los tramos son $3x^2+2ax$ y $b-2/x^2$.
Por tanto la derivabilidad exige $3+2a=b-2$.
Al sustituir $a=b+1$ se obtiene $b=-7$ y $a=-6$.
Estos valores cumplen ambas condiciones.

\textbf{b)} En $x=2$ se usa el segundo tramo. Con $b=3$,
$f(2)=7$ y $f'(2)=3-2/4=5/2$.
La tangente es $y-7=\frac52(x-2)$, o $\boxed{y=\frac52 x+2}$.
